\documentclass[11pt,a4paper]{article}
\usepackage[margin=26mm,headheight=15pt]{geometry}
\usepackage[T1]{fontenc}
\usepackage[utf8]{inputenc}
\usepackage{libertinus}
\usepackage{amsmath,amssymb,amsthm,mathtools}
\usepackage{microtype,booktabs,array,longtable}
\usepackage{xcolor,enumitem,fancyhdr,titlesec,xurl}
\usepackage{aliascnt}
\usepackage{hyperref}
\usepackage[nameinlink,noabbrev]{cleveref}
\definecolor{ink}{RGB}{26,53,74}
\definecolor{muted}{RGB}{90,96,102}
\hypersetup{colorlinks=true,linkcolor=ink,citecolor=ink,urlcolor=ink,
 pdftitle={Localized Spectral Response and the Full Subcritical Pressure Law},
 pdfauthor={Research draft prepared with ChatGPT for Vladimir Reshetnikov},
 pdfsubject={Thue--Morse digital products, moving zeros, weak perturbation, pressure cusps}}
\titleformat{\section}{\Large\bfseries\color{ink}}{\thesection}{.7em}{}
\titleformat{\subsection}{\large\bfseries}{\thesubsection}{.7em}{}
\pagestyle{fancy}\fancyhf{}
\fancyhead[L]{\small Localized spectral response}
\fancyhead[R]{\small ProveIt research}
\fancyfoot[C]{\thepage}
\setlength{\emergencystretch}{3em}
\setlength{\parskip}{.2em}
\setlist{itemsep=.3em,topsep=.5em}
\allowdisplaybreaks
\numberwithin{equation}{section}
\newtheorem{theorem}{Theorem}[section]
\newaliascnt{lemma}{theorem}
\newtheorem{lemma}[lemma]{Lemma}
\aliascntresetthe{lemma}
\newaliascnt{proposition}{theorem}
\newtheorem{proposition}[proposition]{Proposition}
\aliascntresetthe{proposition}
\newaliascnt{corollary}{theorem}
\newtheorem{corollary}[corollary]{Corollary}
\aliascntresetthe{corollary}
\theoremstyle{definition}
\newaliascnt{definition}{theorem}
\newtheorem{definition}[definition]{Definition}
\aliascntresetthe{definition}
\newtheorem{question}{Research question}
\theoremstyle{remark}
\newaliascnt{remark}{theorem}
\newtheorem{remark}[remark]{Remark}
\aliascntresetthe{remark}
\newcommand{\TT}{\mathbb T}
\newcommand{\RR}{\mathbb R}
\newcommand{\CC}{\mathbb C}
\newcommand{\NN}{\mathbb N}
\newcommand{\ZZ}{\mathbb Z}
\newcommand{\dd}{\,\mathrm d}
\newcommand{\Lop}{\mathcal L}
\newcommand{\Id}{\mathrm I}
\newcommand{\PV}{\operatorname{PV}}
\newcommand{\spec}{\operatorname{spec}}
\newcommand{\dist}{\operatorname{dist}}
\newcommand{\norm}[1]{\left\lVert#1\right\rVert}
\newcommand{\abs}[1]{\left|#1\right|}
\newcommand{\eps}{\boldsymbol\varepsilon}
\newcommand{\vel}{\boldsymbol v}
\newcommand{\Proj}{\Pi}
\newcommand{\code}[1]{\texttt{\detokenize{#1}}}
\newcommand{\Holder}{C^\alpha(\TT)}

\begin{document}
\hypersetup{pageanchor=false}
\begin{titlepage}
\thispagestyle{empty}
\setlength{\parindent}{0pt}
{\small\sffamily\color{ink} PROVEIT\quad /\quad ANALYSIS AND DYNAMICAL SYSTEMS}\par
\vspace{17mm}
{\Huge\bfseries\color{ink} Localized Spectral Response\par}
\vspace{5mm}
{\LARGE\bfseries and the Full Subcritical\\[2mm] Pressure Law\par}
\vspace{10mm}
{\Large Moving zeros, a polyhedral cusp,\\[1mm] and finite-size scaling for digital products\par}
\vspace{12mm}
{\large Research draft prepared for Vladimir Reshetnikov\par}
{Prepared with ChatGPT\quad\textbullet\quad 29 September 2026\par}
\vspace{8mm}\hrule\vspace{4mm}
\begin{abstract}
The latest critical-pressure report in the ProveIt repository leaves the
atomic-phase asymptotic unresolved for absolute-moment exponents
$0<s\leq 1/2$. We close that local gap and obtain a quantitative result
throughout $0<s<1$. For the normalized digital mask in every integer
base $b\geq2$,
\[
 P_{b,s}(c)=(1-s)\log b
 +(b-1)\pi s\tan(\pi s/2)|c|
 +O_{b,s,\alpha}(|c|^{1+s-\alpha}),\qquad 0<\alpha<s.
\]
The missing ingredient is a cancellation estimate for the perturbation
paired with the atomic left eigenfunctional. This pairing is $O(|c|)$
on $C^\alpha$, although its total-variation norm is of order
$|c|\log(1/|c|)$ and the full operator changes only by
$O(|c|^{s-\alpha})$.

A stronger multivariate theorem moves the $b-1$ mask zeros independently.
With their displacements denoted by $\eps$, the leading pressure response
is exactly $\pi s\tan(\pi s/2)\norm\eps_1$: the tangent cusp is
polyhedral, not quadratic, and opposite zero motions do not cancel.
We prove a weak distributional expansion, an explicit atomic spectral
gap, the pressure--moment identification, a compact-parameter uniform
version, and an operator finite-size law with an elementary amplitude.
Every repository-specific ingredient is rederived. Reproducible numerical
diagnostics and ten further research questions accompany the proofs.
\end{abstract}
\vfill
{\small\color{muted} Status: unrefereed mathematical research draft. The proofs are not
Lean/Rocq verified. Computations are diagnostics, not interval certificates.
Independent priority and correctness review remain appropriate.}
\end{titlepage}
\hypersetup{pageanchor=true}

\begingroup\small\setlength{\parskip}{0pt}\tableofcontents\endgroup
\newpage

\section{The question being resolved}
\label{sec:scope}

\subsection{A published regularity question and a precise repository gap}
Generalized Thue--Morse products connect digital exponential sums,
Riesz products, and thermodynamic formalism. Gohlke, Kesseb\"ohmer, and
Schindler develop their spectral and fractal theory and obtain an
analytic phase-dependence result at a particular integer pressure order
\cite{GKS}. Gohlke, Lamprinakis, and Schmeling prove continuity in the
singularity location for all nonnegative pressure orders. Following their
Theorem~2.3, they ask about stronger regularity at general positive orders
\cite{GLS}. These statements do not themselves determine a local modulus
or a one-sided derivative at the exceptional atomic phase.

The repository has recently isolated that phase in three reports on
integer, supercritical fractional, and critical pressure. The last of
these, \emph{Critical Cusps in Digital-Product Pressure} \cite{CriticalRepo},
proves a subcritical formula for $1/2<s<1$ and explicitly asks whether it
continues to $0<s\leq1/2$. Its scalar first-order coefficient is already
valid below the cutoff. The obstruction is an error bound of the form
$O(|c|^{2(s-\alpha)})$, obtained by squaring a strong operator norm.
There is no $\alpha>0$ making that error $o(|c|)$ when $s\leq1/2$.

The present article resolves this exact obstruction. The improvement is
not an extrapolation from a first-order calculation. We prove an
$O(|c|)$ bound after pairing the perturbation with the left
eigenfunctional, and combine it with the strong eigenfunction estimate.
The resulting nonlinear error is
\[
 O(|c|)\,O(|c|^{s-\alpha})=O(|c|^{1+s-\alpha})=o(|c|)
\]
for every $0<s<1$ and $0<\alpha<s$.

\subsection{A broader result than a threshold extension}
A common translation moves every zero by the same amount. Instead,
we allow independent small displacements of all the mask zeros while
keeping the modulus of the leading polynomial coefficient fixed. The
main theorem shows that the first pressure increment is a constant
multiple of the sum of the absolute displacements. Thus the method
identifies a multivariate tangent geometry, in addition to removing the
artificial threshold at $s=1/2$.

The proof also distinguishes three different kinds of convergence:
strong operator convergence, convergence after pairing with a regular
test function, and total-variation convergence. Confusing them would
lose exactly the cancellation that makes the result work.

\subsection{What is, and is not, claimed}
\begin{center}
\begin{tabular}{@{}p{.25\textwidth}p{.68\textwidth}@{}}
\toprule
\textbf{Claim class}&\textbf{Scope}\\
\midrule
Proved here & Full fixed-subcritical cusp; quantitative remainder;
independent-zero $\ell^1$ law; weak distributional response; compact-parameter
uniformity; finite-size operator and moment scaling.\\[2mm]
Reproved background & Atomic rank-one spectral gap; translated-zero
coefficient; comparison of polynomial moments with cylinder pressure.\\[2mm]
Not claimed & Global phase regularity; a two-point Lipschitz theorem on
an entire phase neighborhood; a uniform limit as $s\downarrow0$ or
$s\uparrow1$; a sharp next coefficient; formal verification; established
worldwide priority.\\
\bottomrule
\end{tabular}
\end{center}

The repository snapshot used for comparison is
\begin{center}\small
\nolinkurl{458ccfb80b9940220819091da107e82bb970dcca}.
\end{center}
The inspected critical report has Git blob
\nolinkurl{c5491a93464668a8e20dd9afcda664a1c4b6cb81}. Its discussion of the
remaining subcritical interval and its first research question specify
the gap addressed here. The current primary literature and targeted
searches were also checked. This establishes a concrete local
contribution relative to those sources, not a guarantee that no related
formula exists elsewhere. No unverified repository theorem is used as
an axiom in the proofs below.

\section{Definitions and main results}
\label{sec:main}

\subsection{Digital masks with independently moving zeros}
Fix an integer $b\geq2$. Write $T_bx=bx\pmod1$ on
$\TT=\RR/\ZZ$, and let
\[
 z_j=\frac jb,\qquad 1\leq j\leq b-1.
\]
For a real displacement vector
$\eps=(\varepsilon_1,\ldots,\varepsilon_{b-1})$, define
\begin{align}
 p_{\eps}(z)&=\frac1b\prod_{j=1}^{b-1}
       \bigl(z-e^{2\pi i(z_j+\varepsilon_j)}\bigr),
       \label{eq:polynomial}\\
 m_{\eps}(x)&=|p_{\eps}(e^{2\pi i x})|
   =\frac1b\prod_{j=1}^{b-1}2|\sin\pi(x-z_j-\varepsilon_j)|.
       \label{eq:mask}
\end{align}
Only sufficiently small displacements are considered, so the zeros
remain separated. Put $\delta=\norm\eps_\infty$.
At the origin of parameter space,
\begin{equation}
 m_0(x)=\left|\frac1b\sum_{r=0}^{b-1}e^{2\pi i r x}\right|
       =\left|\frac{\sin\pi b x}{b\sin\pi x}\right|,
 \label{eq:atomic-mask}
\end{equation}
with removable values at integers. A common displacement
$\eps=c\mathbf1$ gives $m_{\eps}(x)=m_0(x-c)$ exactly.
For independent displacements, the mask need not be a normalized
probability $g$-function. The digital-product pressure defined below
nevertheless makes sense without that additional property.

For $0<s<1$, the transfer operator is
\begin{equation}
 (\Lop_{s,\eps}f)(x)=\sum_{k=0}^{b-1}
 m_{\eps}\!\left(\frac{x+k}{b}\right)^s
 f\!\left(\frac{x+k}{b}\right).
 \label{eq:operator}
\end{equation}
There is \emph{no factor $1/b$} in this definition.
The operator acts on $\Holder$, with
$\norm f_\alpha=\norm f_\infty+[f]_\alpha$, where the seminorm uses
circle distance. Spectral statements use the complexification.

Let $N=b^n$ and let $\mathcal D_{b,n}$ be the $N$ intervals of length
$1/N$. Define
\begin{align}
 Q_{n,\eps}(z)&=\prod_{r=0}^{n-1}p_{\eps}(z^{b^r}),
 &\deg Q_{n,\eps}&=N-1,\\
 F_{s,n,\eps}(x)&=|Q_{n,\eps}(e^{2\pi i x})|^s,
 &M_{s,n}(\eps)&=\int_0^1F_{s,n,\eps}(x)\dd x,\\
 Z_{s,n}(\eps)&=\sum_{I\in\mathcal D_{b,n}}
       \sup_{x\in I}F_{s,n,\eps}(x),
 &P_{b,s}(\eps)&=\lim_{n\to\infty}\frac1n\log Z_{s,n}(\eps).
 \label{eq:pressure}
\end{align}
Existence of this limit near $\eps=0$ is part of the proof.
For common displacement we abbreviate $P_{b,s}(c)=P_{b,s}(c\mathbf1)$.

The useful atomic quantities are
\begin{equation}
 H(x)=\frac{|\sin\pi x|}{\pi},\quad H_s=H^s,\quad
 \beta_s=b^{1-s},\quad
 \nu_s(f)=\int_0^1\frac{f(x)}{H_s(x)}\dd x,\quad
 \Proj_s=H_s\otimes\nu_s.
 \label{eq:atomic-data}
\end{equation}
Since $s<1$, $\nu_s$ is bounded even on $C^0$, and
$\nu_s(H_s)=1$. Finally set
\begin{equation}
 J_s=\pi s\tan\left(\frac{\pi s}{2}\right),\qquad
 A_s=\frac{2}{\pi s}\tan\left(\frac{\pi s}{2}\right).
 \label{eq:constants}
\end{equation}

\subsection{The full subcritical law and its multivariate strengthening}
\begin{theorem}[Polyhedral pressure cusp]
\label{thm:main}
Let $b\geq2$, $0<s<1$, and $0<\alpha<s$. For all sufficiently small
$\eps$, the pressure exists and equals the logarithm of a real positive,
simple dominant eigenvalue of $\Lop_{s,\eps}$ on $\Holder$. Moreover,
\begin{equation}
 \boxed{\quad
 P_{b,s}(\eps)=(1-s)\log b+J_s\norm\eps_1
                   +O_{b,s,\alpha}(\delta^{1+s-\alpha}).
 \quad}
 \label{eq:main}
\end{equation}
The statement is uniform for $s$ in a compact subset $K\subset(0,1)$
when $0<\alpha<\inf K$ is fixed: the displacement neighborhood and
implicit constant can then be chosen independently of $s\in K$.
\end{theorem}

\begin{corollary}[The previously missing interval]
\label{cor:binary}
For every $0<s<1$,
\begin{equation}
 P_{b,s}(c)=(1-s)\log b+(b-1)J_s|c|
                         +O_{b,s,\alpha}(|c|^{1+s-\alpha}).
 \label{eq:common}
\end{equation}
In the binary convention $\psi_c(x)=\log\cos^2\pi(x-c)$, the pressure
order is $q=s/2$. Thus for all $0<q<1/2$,
\begin{equation}
 P(q\psi_c)=(1-2q)\log2+2\pi q\tan(\pi q)|c|
                     +O_{q,\alpha}(|c|^{1+2q-\alpha}),
 \quad 0<\alpha<2q.
 \label{eq:q-convention}
\end{equation}
In particular, at the formerly excluded endpoint $q=1/4$, the linear
coefficient is $\pi/2$.
\end{corollary}

The sine convention in \cite{GLS} has its atomic phase at $1/2$.
Replacing its parameter by $c+1/2$ gives the cosine convention above;
there is no change in the pressure normalization. In the cosine
convention the classical binary Thue--Morse phase is $c=1/2$, not $c=0$.

\begin{corollary}[Pointwise geometry and one-sided derivatives]
\label{cor:geometry}
The point $\eps=0$ is a strict local minimum of $P_{b,s}$. For every
fixed nonzero real vector $\vel$,
\begin{equation}
 \lim_{t\downarrow0}\frac{P_{b,s}(t\vel)-P_{b,s}(0)}{t}
 =J_s\norm\vel_1.
 \label{eq:direction}
\end{equation}
Consequently the pressure is not differentiable at the atomic point.
For common phase displacement its right and left derivatives at zero
are $(b-1)J_s$ and $-(b-1)J_s$, respectively. It has exact pointwise
H\"older exponent one there.
\end{corollary}
Here pointwise H\"older exponent one means an optimal bound against the
value at the distinguished point. It is not a claim about a uniform
two-point Lipschitz seminorm throughout a neighborhood.

\subsection{Finite-size scaling}
\begin{theorem}[Operator and moment scaling]
\label{thm:finite}
Fix $b,s,\alpha$ as in \cref{thm:main}. There are $C_T<\infty$ and
$\rho\in(0,1)$ such that, whenever $\delta$ is sufficiently small and
$n\delta\leq T$,
\begin{equation}
 \left\|\beta_s^{-n}\Lop_{s,\eps}^{\,n}
       -e^{nJ_s\norm\eps_1}\Proj_s\right\|_{\Holder\to\Holder}
 \leq C_T\bigl(\delta^{s-\alpha}+\rho^n\bigr).
 \label{eq:finite-operator}
\end{equation}
The analogous moment formula is
\begin{equation}
 b^{sn}M_{s,n}(\eps)
   =A_s e^{nJ_s\norm\eps_1}
      +O_T\bigl(\delta^{s-\alpha}+\rho^n\bigr).
 \label{eq:finite-moment}
\end{equation}
In particular, if $\eps=h\vel$, $h\downarrow0$, $n\to\infty$, and
$nh\to\tau\in[0,\infty)$, then
\begin{equation}
 b^{sn}M_{s,n}(h\vel)\longrightarrow
 A_s e^{\tau J_s\norm\vel_1}.
 \label{eq:finite-limit}
\end{equation}
Constants can be made uniform for $s\in K\Subset(0,1)$ and bounded
$T$, with fixed $0<\alpha<\inf K$.
\end{theorem}
The condition $n\to\infty$ in the last statement matters when
$\tau=0$. At fixed $n$ the initial rank-one relaxation has not vanished;
the $\rho^n$ term records that initial layer explicitly.

\section{Why the strong norm alone gives the wrong threshold}
\label{sec:perturbation}

The following elementary perturbation lemma isolates the mechanism. It
is useful whenever a distinguished left eigenfunctional regularizes a
perturbation more effectively than the ambient operator norm.

\begin{lemma}[Strong--weak eigenvalue remainder]
\label{lem:sw}
Let $X$ be a complex Banach space and $L_0\in\mathcal B(X)$. Suppose
$\beta$ is an isolated algebraically simple eigenvalue with
$L_0h=\beta h$, $\nu L_0=\beta\nu$, and $\nu(h)=1$.
For $L=L_0+\Delta$ with $\norm\Delta$ sufficiently small, let $\lambda$
be the nearby eigenvalue. Then
\begin{equation}
 \lambda=\beta+\nu(\Delta h)
       +O\bigl(\norm{\nu\Delta}_{X^*}\norm\Delta_{X\to X}\bigr).
 \label{eq:sw}
\end{equation}
The constant depends on $L_0,h,\nu$ and an isolating contour, but not
on the particular small perturbation.
\end{lemma}
\begin{proof}
Choose a circle $\Gamma$ enclosing only $\beta$. The resolvent identity
and the Neumann series imply that
\[
 \Proj_\Delta=\frac{1}{2\pi i}\int_\Gamma(z-L_0-\Delta)^{-1}\dd z
\]
is defined, has rank one, and satisfies
$\norm{\Proj_\Delta-h\otimes\nu}=O(\norm\Delta)$.
Indeed the inverse on $\Gamma$ is obtained by multiplying the old
resolvent by a uniformly convergent Neumann series. The same contour
formula, or its homotopy in the perturbation parameter, preserves rank.

Normalize
\[
 h_\Delta=\frac{\Proj_\Delta h}{\nu(\Proj_\Delta h)}.
\]
The denominator is $1+O(\norm\Delta)$, hence is nonzero for small
perturbations; $\nu(h_\Delta)=1$ and
$\norm{h_\Delta-h}=O(\norm\Delta)$.
Applying $\nu$ to $Lh_\Delta=\lambda h_\Delta$ gives the exact identity
\[
 \lambda-\beta=\nu(\Delta h_\Delta)
 =\nu(\Delta h)+\nu\Delta(h_\Delta-h).
\]
The last term has the bound in \cref{eq:sw}.
\end{proof}

The customary norm-square estimate follows by replacing
$\norm{\nu\Delta}$ by $\norm\nu\norm\Delta$. Here that replacement
throws away useful information. We will instead show
\begin{equation}
 \norm\Delta_{\Holder\to\Holder}=O(\delta^{s-\alpha}),
 \qquad
 \norm{\nu_s\Delta}_{(\Holder)^*}=O(\delta).
 \label{eq:two-norms}
\end{equation}
Neither estimate alone supplies the theorem. Their product is the
nonlinear remainder that works for the entire subcritical interval.

\section{The atomic operator: an explicit spectral gap}
\label{sec:gap}

\subsection{Coboundary and quadrature}
The atomic mask satisfies, away from removable exceptional points,
\begin{equation}
 m_0(y)^s=\frac{H_s(by)}{b^sH_s(y)}.
 \label{eq:coboundary}
\end{equation}
Multiplication by $H_s$ therefore conjugates the atomic operator to a
uniform inverse-branch average, on the functions for which the quotient
is meaningful. Direct substitution and a change of variables give
\begin{equation}
 \Lop_{s,0}H_s=\beta_sH_s,
 \qquad \nu_s\Lop_{s,0}=\beta_s\nu_s.
 \label{eq:eigenpair}
\end{equation}
The quotient $f/H_s$ can be singular at zero. We do not assume it lies
in an unweighted H\"older space.

\begin{lemma}[Weighted quadrature estimate]
\label{lem:quadrature}
For $0<s<1$ and $0<\alpha<s$, with $N=b^n$,
\begin{equation}
 \left\|\beta_s^{-n}\Lop_{s,0}^{\,n}f
                  -H_s\nu_s(f)\right\|_\infty
 \leq C_{b,s,\alpha}
       \bigl(N^{-\alpha}+N^{s-1}\bigr)\norm f_\alpha.
 \label{eq:quadrature}
\end{equation}
\end{lemma}
\begin{proof}
For $0<x<1$, telescoping \cref{eq:coboundary} gives
\begin{equation}
 \beta_s^{-n}\Lop_{s,0}^{\,n}f(x)
 =H_s(x)\frac1N\sum_{k=0}^{N-1}
       \frac{f((x+k)/N)}{H_s((x+k)/N)}.
 \label{eq:riemann}
\end{equation}
The case $N=1$ is immediate after enlarging the constant. For $N\geq2$,
compare the sum with the integral of $f/H_s$ by splitting off the two
endpoint cells. On these cells $H_s(y)\asymp d(y,\mathbb Z)^s$.
The endpoint integrals are $O(N^{s-1}\norm f_\infty)$ after multiplication
by $H_s(x)$. The first sampled endpoint term is bounded by
\[
 \frac{H_s(x)}{N H_s(x/N)}\norm f_\infty
 \leq C N^{s-1}\norm f_\infty,
\]
and the other endpoint is identical after reflection.

On the interior cells, replacing the sampled numerator by its value at
a point of the cell costs at most
\[
 C N^{-\alpha}[f]_\alpha\int_{1/N}^{1-1/N}H_s(y)^{-1}\dd y
 \leq C N^{-\alpha}\norm f_\alpha.
\]
The variation of the denominator contributes at most a constant times
\[
 N^{-1}\norm f_\infty\int_{1/N}^{1/2}y^{-s-1}\dd y
 \leq C N^{s-1}\norm f_\infty.
\]
Here one may first compare each cell to the minimum and maximum of the
smooth denominator and then sum its variation; no derivative of $f$ is
used. These estimates prove the uniform bound for $0<x<1$.
At $x=0$, only the atomic fixed-point inverse branch survives at every
iteration, and the normalized value is $N^{s-1}f(0)$. Since $H_s(0)=0$,
the same estimate holds there.
\end{proof}

\begin{proposition}[Rank-one exponential relaxation]
\label{prop:gap}
For every $0<s<1$ and $0<\alpha<s$, there exist $C<\infty$ and
$\rho_0<1$ such that
\begin{equation}
 \left\|\beta_s^{-n}\Lop_{s,0}^{\,n}-\Proj_s\right\|_{\Holder\to\Holder}
 \leq C\rho_0^n.
 \label{eq:gap}
\end{equation}
In particular, $\beta_s$ is an algebraically simple isolated dominant
eigenvalue. The constants can be chosen uniformly on compact
$s$-intervals whose lower endpoint exceeds $\alpha$.
\end{proposition}
\begin{proof}
Write $B=\beta_s^{-1}\Lop_{s,0}$. The sup norms of $B^n1$ are uniformly
bounded by \cref{lem:quadrature}. For each fixed $m$, pair corresponding
inverse branches of $T_b^m$; they contract circle distance by $b^{-m}$.
The finite collection of branch weights is $s$-H\"older and hence
$\alpha$-H\"older. The difference of two output values therefore gives
\begin{equation}
 [B^mf]_\alpha\leq C_0 b^{-m\alpha}[f]_\alpha+C_m\norm f_\infty,
 \label{eq:LY}
\end{equation}
where $C_0$ is independent of $m$ and $C_m$ is finite. The first
coefficient comes from the sum of the nonnegative branch weights,
namely $B^m1$; the second comes from their H\"older seminorms.
Across the endpoint of the circle the inverse branches are relabeled
cyclically, giving the same estimate.

The subspace $\ker\nu_s$ is invariant. On it \cref{lem:quadrature}
implies
\[
 \norm{B^nf}_\infty\leq C\theta^n\norm f_\alpha,
 \qquad \theta=\max\{b^{-\alpha},b^{s-1}\}<1.
\]
Choose $m$ with $a=C_0b^{-m\alpha}<1$. If
$u_k=[B^{km}f]_\alpha$, \cref{eq:LY} yields
\[
 u_{k+1}\leq a u_k+C_mC\theta^{km}\norm f_\alpha.
\]
A finite geometric convolution bounds $u_k$ by $C r^k\norm f_\alpha$
for any fixed $r$ strictly between $\max\{a,\theta^m\}$ and one.
The finitely many remaining residues of $n$ modulo $m$ change only the
constant. This proves exponential contraction on $\ker\nu_s$ in the
full norm. Since $\Proj_s$ is a commuting rank-one projection by
\cref{eq:eigenpair}, \cref{eq:gap} follows.

An additional eigenvalue of modulus at least one for $B$, or a
nontrivial Jordan vector at one, contradicts this exponential convergence.
The spectrum on the kernel is contained in a smaller disk by the
spectral radius formula, proving the stated spectral conclusions.
All estimates are uniform when $s$ stays away from zero and one and
$\alpha$ is fixed below its lower endpoint; the same $m$ may then be used.
\end{proof}

\section{Translated zeros as weak distributions}
\label{sec:distribution}

\subsection{The universal one-zero constant}
The following integral is the local contribution of opening a simple
zero. Its evaluation appears in the predecessor report; we include the
proof to keep the new argument independent.

\begin{lemma}[Model singular integral]
\label{lem:J}
For $0<s<1$,
\begin{equation}
 \int_0^\infty u^{-s}
       \bigl((u+1)^s+|u-1|^s-2u^s\bigr)\dd u
 =J_s=\pi s\tan(\pi s/2).
 \label{eq:J}
\end{equation}
The paired integrand is integrable at both zero and infinity.
\end{lemma}
\begin{proof}
At zero it is $O(u^{-s})$; at infinity the two linear Taylor terms
cancel, leaving $O(u^{-2})$. Substitute $u=1/v$ and integrate by parts.
The endpoint terms vanish and the result, divided by $s$, is
\begin{equation}
 \int_0^\infty
 \frac{(1+v)^{s-1}+\operatorname{sgn}(v-1)|v-1|^{s-1}}{v}\dd v.
 \label{eq:Jparts}
\end{equation}
The terms at zero must remain paired. Transform the first term using
$t=v/(1+v)$ and pair it with the negative term on $(0,1)$. Their sum is
\[
 \int_0^1\frac{(1-t)^{-s}-(1-t)^{s-1}}t\dd t
 =\psi(s)-\psi(1-s)=-\pi\cot(\pi s).
\]
The remaining term is
\[
 \int_1^\infty\frac{(v-1)^{s-1}}v\dd v
 =B(s,1-s)=\pi\csc(\pi s).
\]
The beta and digamma reflection identities used here are standard
\cite[\S5.5]{DLMF}. Their sum is
$\pi(\csc\pi s-\cot\pi s)=\pi\tan(\pi s/2)$.
\end{proof}

\subsection{A quantitative distributional expansion}
For a point $z\in\TT$, define
\[
 q_{z,h}(x)=\left|\frac{\sin\pi(x-z-h)}{\sin\pi(x-z)}\right|^s-1,
 \qquad k_z(x)=\pi\cot\pi(x-z).
\]
Values at the zero denominator are irrelevant for integration. The
notation $\PV$ always means deletion of a symmetric interval around
that zero, followed by a limit.

\begin{lemma}[One-zero weak derivative]
\label{lem:onezero}
Let $0<s<1$ and $0<\alpha<1$. For $g\in C^\alpha(\TT)$,
\begin{equation}
 \int_0^1q_{z,h}(x)g(x)\dd x
 =J_s|h|g(z)-sh\PV\!\int_0^1 k_z(x)g(x)\dd x
     +O_{s,\alpha}(|h|^{1+\alpha}\norm g_\alpha).
 \label{eq:onezero}
\end{equation}
The principal-value functional is bounded on $C^\alpha$. In particular,
the left side is $O(|h|\norm g_\alpha)$.
For $g=1$ the error can be improved to $O_s(h^2)$.
\end{lemma}
\begin{proof}
Translation lets us take $z=0$. Choose a fixed small symmetric interval
$(-r,r)$ about zero and use its signed real coordinate. There,
$|\sin\pi x|=|x|a(x)$ with $a$ positive and smooth. Set
\[
 U_h(x)=|x-h|^s|x|^{-s},\qquad A(x)=a'(x)/a(x).
\]
Uniformly on this fixed interval,
\begin{equation}
 q_{0,h}(x)=U_h(x)-1-shA(x)U_h(x)+O(h^2U_h(x)).
 \label{eq:local-factor}
\end{equation}
The integrals of $U_h$ on $(-r,r)$ are uniformly bounded.

First take a constant test. Pairing positive and negative $x$ and
scaling by $|h|$ gives
\[
 \int_{-r}^{r}(U_h-1)\dd x
 =|h|J_s+O(h^2),
\]
because the tail of the paired integral in \cref{eq:J} is $O(1/R)$.
The smooth correction in \cref{eq:local-factor} contributes
$-sh\int_{-r}^{r}A+O(h^2)$. To see the stated error, use the estimate
$\int(U_h-1)A=O(|h|)$: subtract $A(0)$, use its Lipschitz modulus,
and pair the constant part. On the complement, ordinary Taylor
expansion contributes $-sh\int k_0+O(h^2)$.
The local and complementary linear terms add to
$-sh\PV\int k_0=0$ by periodicity. This proves the constant-test
assertion.

For a general test, subtract $g(0)$ in the local interval. On
$|x|\leq2|h|$, its modulus is at most $[g]_\alpha|x|^\alpha$.
The integrals of
\[
 (U_h+1)|x|^\alpha
 \quad\hbox{and}\quad |h||x|^{\alpha-1}
\]
over that region are $O(|h|^{1+\alpha})$; integrability of the first
uses $s<1$ (indeed $\alpha-s>-1$). On
$2|h|<|x|<r$, Taylor expansion gives
\[
 U_h-1+\frac{sh}{x}=O(h^2/x^2).
\]
Multiplying by $g(x)-g(0)$ and integrating costs
\[
 C h^2[g]_\alpha\int_{2|h|}^{r}x^{\alpha-2}\dd x
 \leq C|h|^{1+\alpha}\norm g_\alpha.
\]
The smooth factors in \cref{eq:local-factor} and the complement have
the same or a smaller remainder. One may use the elementary estimate
$\int_{-r}^{r}|U_h-1|\dd x=O(|h|\log(1/|h|))$; multiplying by
$|h|$ is still $O(|h|^{1+\alpha})$ since $\alpha<1$.

Finally,
$\PV\int_{-r}^{r}g(x)/x=\int_{-r}^{r}(g(x)-g(0))/x$ is absolutely
convergent and bounded by $C[g]_\alpha$. The remaining smooth part is
bounded by $C\norm g_\infty$. This proves both the expansion and the
boundedness of its principal-value term.
\end{proof}

In distribution notation, for $h\downarrow0$,
\begin{equation}
 h^{-1}q_{z,\pm h}(x)\dd x
 \longrightarrow J_s\delta_z\mp s\PV(k_z(x)\dd x)
 \quad\hbox{in }(C^\alpha)^*.
 \label{eq:distribution}
\end{equation}
The Dirac term is even in the displacement; the principal-value term
is odd. The Dirac term cannot be recovered by pointwise differentiation
away from the zero.

\subsection{Several zeros and locally regular test functions}
Let
\begin{equation}
 R_{\eps}(x)=\left(\frac{m_{\eps}(x)}{m_0(x)}\right)^s
          =\prod_{j=1}^{b-1}(1+q_{z_j,\varepsilon_j}(x)).
 \label{eq:ratio}
\end{equation}
Fix disjoint closed neighborhoods $U_j$ of the $z_j$, none containing
zero. It is useful to allow a test function that is only integrable
away from these neighborhoods. Write
\[
 \norm g_* =\norm g_{L^1(\TT)}+
       \sum_{j=1}^{b-1}\norm{g|_{U_j}}_{C^\alpha(U_j)}.
\]
Shrinking a neighborhood slightly when multiplying by a fixed cutoff
only changes constants.

\begin{proposition}[Multizero weak expansion]
\label{prop:multizero}
For $0<s<1$, $0<\alpha<1$, and tests with finite $\norm g_*$,
\begin{align}
 \int_0^1(R_{\eps}-1)g\dd x
 &=J_s\sum_{j=1}^{b-1}|\varepsilon_j|g(z_j)
   -s\sum_{j=1}^{b-1}\varepsilon_j\PV\!\int_0^1 k_{z_j}(x)g(x)\dd x
   +O(\delta^{1+\alpha}\norm g_*).
 \label{eq:multizero}
\end{align}
In particular this is $O(\delta\norm g_*)$. For $g=1$,
\begin{equation}
 \int_0^1(R_{\eps}-1)\dd x
       =J_s\norm\eps_1+O(\delta^2).
 \label{eq:scalar}
\end{equation}
Constants are uniform when $s$ ranges over a compact subset of $(0,1)$.
\end{proposition}
\begin{proof}
Use a smooth partition of unity subordinate to the $U_j$ and their
complement. On a neighborhood of $z_j$, every factor with index
$k\ne j$ has a smooth expansion
\[
 q_{z_k,\varepsilon_k}=-s\varepsilon_k k_{z_k}+O(\delta^2)
\]
with uniform bounds on a fixed number of derivatives. On the
complement all factors have such expansions, and the error is bounded
pointwise by $C\delta^2$, which can be integrated against $|g|$.

For the terms linear in one $q_j$, apply \cref{lem:onezero} locally,
using cutoffs, and use Taylor expansion elsewhere. A principal value
at $z_j$ is well defined because $g$ is H\"older there; elsewhere its
kernel is bounded against the integrable test.

It remains to control mixed products of two or more distinct $q_j$.
Near $z_j$, a mixed term containing $q_j$ has the form $q_j a g$,
where $a$ is smooth with $C^1$ norm $O(\delta)$, or smaller. Apply the
$O(|\varepsilon_j|)$ functional bound from \cref{lem:onezero} to the
cutoff test $ag$. The result is $O(\delta^2\norm g_*)$.
A mixed term not containing that local singular factor is pointwise
$O(\delta^2)$. Thus all mixed terms are $O(\delta^2\norm g_*)$,
not $O(\delta^2\log(1/\delta))$. This signed estimate is the place
where absolute-value estimates would be needlessly wasteful.

For $g=1$, each unmixed term has the improved $O(\varepsilon_j^2)$
remainder from \cref{lem:onezero}, and
$\PV\int_0^1 k_{z_j}=0$. Mixed terms retain the $O(\delta^2)$ bound.
This proves \cref{eq:scalar}. Separation of the fixed zeros and compactness
of the exponent interval give the stated uniformity.
\end{proof}

\section{The localized left-eigenfunctional estimate}
\label{sec:weak}

\subsection{An exact change-of-variables identity}
Let $\Delta_{\eps}=\Lop_{s,\eps}-\Lop_{s,0}$. Combining the
change of variables $x=T_by$ with \cref{eq:coboundary} gives
\begin{equation}
 \nu_s(\Delta_{\eps} f)
 =\beta_s\int_0^1(R_{\eps}(y)-1)\frac{f(y)}{H_s(y)}\dd y.
 \label{eq:weak-exact}
\end{equation}
This identity is valid for every continuous $f$. Near a mask zero, the
ratio has a $|y-z_j|^{-s}$ singularity; near the atomic point zero,
$H_s^{-1}$ is integrable and the ratio is smooth. Hence all its
integrals exist when $s<1$.

Define bounded functionals on $\Holder$ by
\begin{equation}
 E_j(f)=\frac{f(z_j)}{H_s(z_j)},\qquad
 D_j(f)=\PV\!\int_0^1 k_{z_j}(y)\frac{f(y)}{H_s(y)}\dd y.
 \label{eq:functionals}
\end{equation}
Near $z_j$ the denominator is smooth and nonzero. Near zero the
cotangent kernel is bounded, so the remaining integrability follows
from $s<1$. Thus $D_j$ is well defined even though $f/H_s$ need not
be a globally continuous function.

\begin{theorem}[Weak operator expansion]
\label{thm:weak}
For $0<s<1$ and $0<\alpha<s$,
\begin{equation}
 \nu_s\Delta_{\eps}
 =\beta_s\left(J_s\sum_j|\varepsilon_j|E_j
                   -s\sum_j\varepsilon_jD_j\right)
       +O_{(\Holder)^*}(\delta^{1+\alpha}).
 \label{eq:weak-expansion}
\end{equation}
Consequently,
\begin{align}
 \norm{\nu_s\Delta_{\eps}}_{(\Holder)^*}&\leq C\delta,
 \label{eq:weak-bound}\\
 \nu_s(\Delta_{\eps} H_s)&=\beta_sJ_s\norm\eps_1+O(\delta^2).
 \label{eq:weak-scalar}
\end{align}
\end{theorem}
\begin{proof}
For $g=f/H_s$ in \cref{prop:multizero},
\[
 \norm g_{L^1}\leq\norm f_\infty\int H_s^{-1},\qquad
 \norm{g|_{U_j}}_{C^\alpha}\leq C\norm f_\alpha.
\]
Thus $\norm g_*\leq C\norm f_\alpha$, and \cref{eq:multizero} applied
to \cref{eq:weak-exact} proves \cref{eq:weak-expansion}.
For $f=H_s$ the quotient test is exactly one. Therefore every $E_j$
equals one, every $D_j$ equals zero, and the improved scalar estimate
\cref{eq:scalar} proves \cref{eq:weak-scalar}.
\end{proof}

\subsection{Why a total-variation bound cannot replace cancellation}
\begin{proposition}[Sharp logarithmic loss in the continuous dual norm]
\label{prop:TV}
For common displacement $\eps=c\mathbf1$ and fixed $0<s<1$,
\begin{equation}
 \norm{\nu_s\Delta_{c\mathbf1}}_{(C^0)^*}
 =2\beta_s s\left(\sum_{j=1}^{b-1}H_s(z_j)^{-1}\right)
       |c|\log\frac1{|c|}+O(|c|).
 \label{eq:TV}
\end{equation}
In particular, an $O(|c|)$ estimate in this larger dual norm is false.
\end{proposition}
\begin{proof}
By \cref{eq:weak-exact}, the functional is represented by the signed
integrable density $\beta_s(R_{c\mathbf1}-1)H_s^{-1}$. Its norm on
$C^0$ is its total variation; continuous tests approximate its sign
in the corresponding $L^1$ space.
In a signed coordinate $u$ around $z_j$, for
$2|c|<|u|<r$,
\[
 R_{c\mathbf1}(z_j+u)-1=-\frac{sc}{u}
                  +O(c^2/u^2)+O(|c|).
\]
Also $H_s(z_j+u)^{-1}=H_s(z_j)^{-1}+O(|u|)$.
The two sides of the puncture therefore contribute
$2sH_s(z_j)^{-1}|c|\log(1/|c|)+O(|c|)$.
The inner interval $|u|\leq2|c|$ contributes $O(|c|)$ by scaling and
$s<1$. Away from all the zeros the ratio difference is $O(|c|)$,
and the atomic weight remains integrable. Summing and multiplying by
$\beta_s$ gives \cref{eq:TV}.
\end{proof}

Thus the better estimate is genuinely weak: a H\"older test cannot
follow the rapidly changing sign around a translated zero at unit
amplitude. Pairing the two sides cancels the logarithmic divergence.
The atom in \cref{eq:distribution} survives that cancellation and
produces the cusp coefficient.

\section{Strong perturbation and the pressure theorem}
\label{sec:proof-main}

\subsection{Translation in H\"older norm}
\begin{lemma}[Strong operator modulus]
\label{lem:strong}
For $0<s<1$, $0<\alpha<s$, and sufficiently small $\eps$,
\begin{equation}
 \norm{\Lop_{s,\eps}-\Lop_{s,0}}_{\Holder\to\Holder}
 \leq C\delta^{s-\alpha}.
 \label{eq:strong}
\end{equation}
\end{lemma}
\begin{proof}
Each function $u_j(x)=|\sin\pi(x-z_j)|^s$ is $s$-H\"older. Its
translation difference satisfies
\[
 \norm{u_j(\cdot-h)-u_j}_\infty\leq C|h|^s,
 \qquad [u_j(\cdot-h)-u_j]_s\leq C.
\]
Splitting the H\"older quotient at distance $|h|$, or using the
minimum of these two estimates, gives
\[
 \norm{u_j(\cdot-h)-u_j}_{C^\alpha}\leq C|h|^{s-\alpha}.
\]
The identity
$m_{\eps}^s=b^{-s}2^{s(b-1)}\prod_j u_j(\cdot-\varepsilon_j)$,
a telescoping product, and the Banach algebra estimate for $C^\alpha$
show that the weight difference has this same bound with
$|h|$ replaced by $\delta$. The finite sum of contracted inverse-branch
compositions in \cref{eq:operator} is bounded on $C^\alpha$.
Cyclic relabeling at the circle endpoint again makes this a global
estimate.
\end{proof}

\subsection{The eigenvalue expansion}
\begin{proposition}[Quantitative spectral response]
\label{prop:eigenvalue}
Under the hypotheses of \cref{thm:main}, the eigenvalue near $\beta_s$
is real, positive, and satisfies
\begin{equation}
 \lambda_s(\eps)=\beta_s\left(1+J_s\norm\eps_1\right)
                       +O(\delta^{1+s-\alpha}).
 \label{eq:eigenvalue}
\end{equation}
It is simple and strictly larger in modulus than all other spectrum.
If $\Proj_{s,\eps}$ is its spectral projection, then
\begin{equation}
 \norm{\Proj_{s,\eps}-\Proj_s}=O(\delta^{s-\alpha}).
 \label{eq:projection}
\end{equation}
\end{proposition}
\begin{proof}
Apply \cref{lem:sw} using \cref{prop:gap,thm:weak,lem:strong}. The
linear scalar term is \cref{eq:weak-scalar}; the nonlinear error is
$O(\delta^{1+s-\alpha})$. The scalar $O(\delta^2)$ remainder is smaller,
since $s-\alpha<1$. This proves \cref{eq:eigenvalue} and the contour
construction gives \cref{eq:projection}.

The atomic spectral gap supplies a fixed separating annulus. A
Neumann-series resolvent argument on its compact boundary, together
with the usual large-$z$ resolvent bound, keeps the complementary
spectrum in a disk of radius strictly below $\beta_s$ for all small
perturbations. The nearby eigenvalue tends to $\beta_s$ and stays
outside that disk. Complex conjugation preserves the operator and the
isolating circle, so uniqueness of the eigenvalue in the circle makes
it real; it is positive by its limit. These conclusions do not rely
on a positivity theorem for a strictly positive potential.
\end{proof}

\subsection{Why cylinder pressure equals moment growth}
This step is needed because a singular logarithmic potential should
not be placed into a classical positive-potential theorem without
checking its hypotheses.

\begin{lemma}[Polynomial sampling bound]
\label{lem:sampling}
For every fixed $s>0$ there is $C_s<\infty$ such that, for every
polynomial $p$ of degree at most $N-1$,
\begin{equation}
 N\int_0^1|p(e^{2\pi i x})|^s\dd x
 \leq\sum_{I\in\mathcal D_N}\sup_{x\in I}|p(e^{2\pi i x})|^s
 \leq C_s N\int_0^1|p(e^{2\pi i x})|^s\dd x.
 \label{eq:sampling}
\end{equation}
Here $\mathcal D_N$ is the uniform partition into $N$ intervals.
\end{lemma}
\begin{proof}
The lower bound is immediate by integration on each cell. For the
upper bound choose $R=1+1/N$. Subharmonicity of $|p|^s$ and the Poisson
majorant on the disk of radius $R$ give
\[
 |p(e^{2\pi i x})|^s\leq
 \int_0^1P_{1/R}(2\pi(x-y))|p(Re^{2\pi i y})|^s\dd y.
\]
The Poisson kernel is bounded by
$C N/(1+N^2\dist(x-y,\mathbb Z)^2)$. For fixed $y$, summing its
supremum over the $N$ cells gives at most $CN$: the bounded number
of nearest cells contribute $O(N)$ and the others contribute a
convergent $N\sum_{k\geq1}k^{-2}$ tail. Consequently the cylinder sum
is bounded by
$CN\int_0^1|p(Re^{2\pi i y})|^s\dd y$.

If $D=\deg p$, the reciprocal polynomial $q(z)=z^Dp(1/z)$ and the
monotonicity of radial means of a subharmonic function show that
\[
 \int_0^1|p(Re^{2\pi i y})|^s\dd y
 \leq R^{Ds}\int_0^1|p(e^{2\pi i y})|^s\dd y.
\]
Indeed evaluate the radial mean of $|q|^s$ at radius $1/R$ and compare
it with radius one. Since $D\leq N-1$, $R^{Ds}\leq e^s$.
This proves the upper bound. Constant and zero polynomials are covered
by the same inequality directly.
\end{proof}

A change of variables in the transfer operator gives the exact identity
\begin{equation}
 \int_0^1\Lop_{s,\eps}^{\,n}1\dd x=b^nM_{s,n}(\eps).
 \label{eq:moment-transfer}
\end{equation}
The spectral projection amplitude at the atomic point is
\begin{equation}
 \int_0^1\Proj_s1\dd x
 =\left(\int_0^1H_s\right)\left(\int_0^1H_s^{-1}\right)>0.
 \label{eq:amp-positive}
\end{equation}
By \cref{eq:projection} this amplitude remains positive for small real
$\eps$. The spectral decomposition therefore gives the logarithmic
growth rate $\log\lambda_s(\eps)$ to the left side of
\cref{eq:moment-transfer}. Applying \cref{lem:sampling} to
$Q_{n,\eps}$ yields
\begin{equation}
 b^nM_{s,n}(\eps)\leq Z_{s,n}(\eps)
                     \leq C_s b^nM_{s,n}(\eps).
 \label{eq:pressure-comparison}
\end{equation}
Thus the pressure limit exists and equals $\log\lambda_s(\eps)$.

\begin{proof}[Completion of \cref{thm:main}]
Take the logarithm of \cref{eq:eigenvalue}. Its quadratic logarithmic
correction is $O(\delta^2)$ and is absorbed into
$O(\delta^{1+s-\alpha})$. The preceding comparison identifies the
result with the pressure in \cref{eq:pressure}. Compact-parameter
uniformity follows from the uniform estimates in the spectral gap,
the local integral lemmas, and the strong translation bound. The
isolating resolvent circles may be chosen at a fixed relative distance
from $\beta_s$ because the normalized gap is uniform on such compact
sets. This proves every assertion of the theorem.
\end{proof}

\begin{proof}[Proof of \cref{cor:geometry}]
Since $J_s>0$ and $\norm\eps_1\geq\delta$, the error in
\cref{eq:main} is smaller than $(J_s/2)\norm\eps_1$ for sufficiently
small nonzero $\eps$. This proves the strict local minimum. Substitute
$\eps=t\vel$ and divide by $t>0$ to obtain \cref{eq:direction}.
The directional derivative is even in $\vel$ and positive for nonzero
$\vel$, so it cannot be a linear functional. Common displacement gives
the displayed left and right derivatives and the optimal pointwise
exponent.
\end{proof}

\section{Finite-size law and explicit amplitude}
\label{sec:finite-proof}

The contour construction gives a uniform decomposition
\begin{equation}
 \beta_s^{-n}\Lop_{s,\eps}^{\,n}
 =\left(\frac{\lambda_s(\eps)}{\beta_s}\right)^n\Proj_{s,\eps}
       +R_{n,s,\eps},\qquad
 \norm{R_{n,s,\eps}}\leq C\rho^n
 \label{eq:spectral-powers}
\end{equation}
for some $\rho<1$. For example, integrate the complementary resolvent
on a circle in the normalized spectral gap. Its uniformly bounded
resolvent gives the displayed estimate, with a slight enlargement of
the spectral radius if necessary.

\begin{proof}[Proof of \cref{thm:finite}]
By \cref{eq:eigenvalue},
\[
 \log\frac{\lambda_s(\eps)}{\beta_s}
 =J_s\norm\eps_1+O(\delta^{1+s-\alpha}).
\]
When $n\delta\leq T$, the accumulated logarithmic error is bounded
by $C_T\delta^{s-\alpha}$. Both the main exponential and the perturbed
exponential are uniformly bounded on this time window. Therefore
\[
 \left|\left(\frac{\lambda_s(\eps)}{\beta_s}\right)^n
             -e^{nJ_s\norm\eps_1}\right|
 \leq C_T\delta^{s-\alpha}.
\]
Combine this with \cref{eq:projection,eq:spectral-powers} to obtain
\cref{eq:finite-operator}.

Apply the operator to $1$ and integrate. By
\cref{eq:moment-transfer}, the resulting normalization is
$\beta_s^{-n}b^n=b^{sn}$. The beta-integral evaluation
\[
 \int_0^1|\sin\pi x|^a\dd x
 =\frac{\Gamma((a+1)/2)}{\sqrt\pi\,\Gamma(1+a/2)},\qquad a>-1,
\]
and the reflection identities \cite{DLMF} give
\begin{align}
 \left(\int H_s\right)\left(\int H_s^{-1}\right)
 &=\frac{\Gamma((1+s)/2)\Gamma((1-s)/2)}
          {\pi\Gamma(1+s/2)\Gamma(1-s/2)}\\
 &=\frac{2}{\pi s}\tan(\pi s/2)=A_s.
 \label{eq:amplitude}
\end{align}
This proves \cref{eq:finite-moment}. With $\eps=h\vel$,
$n\norm\eps_1=nh\norm\vel_1$ and the two error terms tend to zero
under the hypotheses of \cref{eq:finite-limit}.
The same argument uses uniform constants on the stated compact
parameter sets.
\end{proof}

An exact relation between the cusp and the amplitude is
\begin{equation}
 J_s=\frac{\pi^2s^2}{2}\,A_s.
 \label{eq:amplitude-cusp}
\end{equation}
The factor $b-1$ in a common translation counts zeros; the amplitude
$A_s$ itself is independent of the base. Neither fact extends to an
arbitrary changed atomic mask without further analysis.

\section{Consequences, examples, and normalization boundaries}
\label{sec:examples}

\subsection{Two concrete values in the newly covered range}
At $s=1/4$, corresponding to pressure order $q=1/8$ in the binary
squared-mask convention,
\[
 J_{1/4}=\frac\pi4(\sqrt2-1),\qquad
 P_{2,1/4}(c)=\frac34\log2+\frac\pi4(\sqrt2-1)|c|
           +O_\alpha(|c|^{5/4-\alpha}).
\]
At $s=1/2$, corresponding to $q=1/4$,
\[
 J_{1/2}=\frac\pi2,\qquad
 P_{2,1/2}(c)=\frac12\log2+\frac\pi2|c|
           +O_\alpha(|c|^{3/2-\alpha}).
\]
The second point is not a new phase transition. It was only the point
where the old norm-square proof ceased to control its remainder.

\subsection{Opposite zero motions do not cancel}
For $b=3$, move the zeros at $1/3$ and $2/3$ by $t$ and $-t$.
Their mean displacement is zero, but
\begin{equation}
 P_{3,s}(t,-t)=(1-s)\log3+2J_s|t|
                         +O(|t|^{1+s-\alpha}).
 \label{eq:opposite}
\end{equation}
A perturbation rule based only on the center of the zero configuration
would incorrectly predict no first-order response. The pressure sees
the total absolute opening of the simple zeros.

\subsection{Smooth zero trajectories and amplitude changes}
\begin{corollary}[Trajectory law]
\label{cor:trajectory}
Suppose $\varepsilon_j(t)=v_jt+O(t^2)$ as $t\to0$. For fixed
$0<s<1$ and $0<\alpha<s$,
\begin{equation}
 P_{b,s}(\eps(t))=(1-s)\log b+J_s\sum_j|v_j|\,|t|
                              +O(|t|^{1+s-\alpha}).
 \label{eq:trajectory}
\end{equation}
If the mask is additionally multiplied by a positive factor $e^{a(t)}$,
its pressure is increased by exactly $sa(t)$.
\end{corollary}
\begin{proof}
The absolute value is Lipschitz, so
$\norm{\eps(t)}_1=|t|\sum|v_j|+O(t^2)$; apply \cref{thm:main}.
Multiplication of the mask by $e^{a(t)}$ multiplies every transfer
operator by $e^{sa(t)}$, or every $n$-step product by $e^{nsa(t)}$.
Taking logarithmic growth gives the exact second assertion.
\end{proof}
This normalization clause is essential. A varying leading coefficient
adds an ordinary signed linear term and can obscure the cusp. The
$\ell^1$ law is stated for the fixed coefficient modulus in
\cref{eq:polynomial}.

\subsection{Digital coefficients and the Thue--Morse connection}
Write $p_{\eps}(z)=\sum_{r=0}^{b-1}a_r(\eps)z^r$. If
$k=\sum_{j=0}^{n-1}d_jb^j$ with $0\leq d_j<b$, then
\[
 [z^k]Q_{n,\eps}(z)=\prod_{j=0}^{n-1}a_{d_j}(\eps).
\]
Thus independently moving the zeros produces a natural family of
complex digit weights. It does not merely perturb an abstract
transfer operator.

For common phase displacement one may choose the equal-modulus
polynomial with coefficients $b^{-1}e^{-2\pi i cr}$. Its product is
\[
 b^{-n}\sum_{k=0}^{b^n-1}e^{-2\pi i c\,s_b(k)}e^{2\pi i kx},
\]
where $s_b(k)$ is the base-$b$ digit sum. Its modulus is the one used
throughout this article. For $b=2,c=1/2$, the coefficients reduce to
$(-1)^{s_2(k)}$, the Thue--Morse signs. At the atomic phase $c=0$ the
polynomial is the normalized Dirichlet polynomial
$b^{-n}\sum_{k=0}^{b^n-1}z^k$. The singular limit studied here therefore
starts from the simplest digital product but has nontrivial fractional
moment response.

\subsection{Endpoint information that does not justify exchanging limits}
The explicit constants satisfy
\begin{align}
 J_s&=\frac{\pi^2s^2}{2}+\frac{\pi^4s^4}{24}+O(s^6),
 &A_s&=1+\frac{\pi^2s^2}{12}+O(s^4)
 &&(s\downarrow0),\\
 J_s&\sim\frac{2}{1-s},
 &A_s&\sim\frac{4}{\pi^2(1-s)}
 &&(s\uparrow1).
 \label{eq:endpoint-constants}
\end{align}
These are expansions of the explicit constants, not uniform
joint-limit theorems for pressure. The proof uses a H\"older exponent
below $s$, and the dual density ceases to be integrable at $s=1$.
At criticality the predecessor report finds a different spectral
mechanism involving a two-dimensional collision \cite{CriticalRepo}.
The present theorem does not replace that critical analysis.

\section{Numerical diagnostics and reproducibility}
\label{sec:numerics}

\subsection{What was actually run}
The package includes \code{code/verify.py}, its complete output
\code{data/verification.json}, and \code{data/run.log}. The included run
used the \code{--full} option. Its checks have three independent roles:
to evaluate the universal singular integral, to test the signed weak
response on a nonconstant function, and to compare finite-dimensional
transfer approximations with the pressure and finite-size formulas.
They do not form part of the analytic proofs.

The $J_s$ integral was evaluated at 65 decimal digits after substitutions
removing the algebraic singularity at zero and mapping the infinite
tail to a finite interval. At $s=0.1,0.25,0.5,0.75,0.9$ its discrepancy
from the elementary expression was below $10^{-30}$. This is a
high-precision quadrature check, not a rigorous interval enclosure.

\subsection{A test that detects the principal-value sign}
For the single binary zero $z=1/2$, use
\[
 g(x)=1+0.2\cos(2\pi x)+0.1\sin(2\pi x).
\]
Then $g(z)=0.8$ and
\[
 \PV\int_0^1\pi\cot\pi(x-1/2)g(x)\dd x=-0.1\pi.
\]
The sine term follows from
$\tan(\pi x)\sin(2\pi x)=2\sin^2(\pi x)$; the other terms have
zero principal value by reflection. Hence the limits of
$|h|^{-1}\int q_{z,\pm h}g$ are
\[
 0.8J_s\pm0.1\pi s.
\]
The code checks both signs at four exponents and three phase sizes.
A check with the constant test alone would miss an incorrect sign in
the principal-value part.

The double-precision adaptive quadrature emitted a roundoff convergence
warning in this diagnostic family. That warning is retained in the run
log. The output includes the quadrature error estimates; it is not
silently promoted to a certified result. The integral identity and its
sign are established by \cref{lem:onezero}, not by this computation.

\subsection{Pressure: mesh comparisons}
Collocation uses a uniform periodic grid and linear interpolation at
each inverse-branch point. The resulting operator is nonnegative.
Power iteration is stopped when the sup norm change of the normalized
iterate is less than $3\times10^{-13}$; the matrix residual is recorded.
The coarse and fine meshes below have $262\,144$ and $524\,288$ points.
The table shows
\[
 \frac{P^{\mathrm{grid}}_{b,s}(h\vel)-(1-s)\log b}{h}.
\]
For $b=2$ the velocity is $(1)$; for $b=3$ it is $(1,-1)$.
\begin{center}\small
\begin{tabular}{@{}rrrrrr@{}}
\toprule
$b$ & $s$ & $h$ & coarse grid & fine grid & predicted limit\\
\midrule
2 & 0.20 & 0.01 & 0.180110 & 0.180113 & 0.204153\\
2 & 0.20 & 0.003 & 0.194370 & 0.194378 & 0.204153\\
2 & 0.20 & 0.001 & 0.200329 & 0.200320 & 0.204153\\
2 & 0.25 & 0.01 & 0.284731 & 0.284733 & 0.325323\\
2 & 0.25 & 0.003 & 0.308837 & 0.308843 & 0.325323\\
2 & 0.25 & 0.001 & 0.318789 & 0.318782 & 0.325323\\
2 & 0.50 & 0.01 & 1.273868 & 1.273868 & 1.570796\\
2 & 0.50 & 0.003 & 1.446803 & 1.446805 & 1.570796\\
2 & 0.50 & 0.001 & 1.519991 & 1.519989 & 1.570796\\
3 & 0.30 & 0.01 & 0.952195 & 0.952194 & 0.960433\\
3 & 0.30 & 0.003 & 0.958221 & 0.958219 & 0.960433\\
3 & 0.30 & 0.001 & 0.959884 & 0.959841 & 0.960433\\
\bottomrule
\end{tabular}

\end{center}
The low-exponent binary rows approach the predicted coefficient, and
the base-three opposite-motion rows distinguish the $\ell^1$ law from
a mean-displacement rule. Mesh changes remain visible after division
by $h$. A small matrix residual only checks the discretized eigenproblem;
it does not bound the distance to the infinite-dimensional spectrum.
No number in the table is labeled a certified continuum digit.

\subsection{Finite-size amplitude}
For $b=2,s=1/4,\vel=(1)$ and $nh=0.1$, the prediction is
$A_{1/4}e^{0.1J_{1/4}}$. Direct iteration of the normalized collocation
operator on the fine mesh gives:
\begin{center}\small
\begin{tabular}{@{}rrrr@{}}
\toprule
$n$ & $h$ & $b^{sn}M_n$ (grid) & limiting value\\
\midrule
25 & 0.00400 & 1.0780086 & 1.0896650\\
50 & 0.00200 & 1.0828660 & 1.0896650\\
100 & 0.00100 & 1.0857978 & 1.0896650\\
200 & 0.00050 & 1.0875046 & 1.0896650\\
\bottomrule
\end{tabular}

\end{center}
The convergence includes both phase error and finite-size relaxation;
it is not expected to be exact at a fixed mesh or a fixed number of
steps. The gamma-product and elementary expressions for $A_s$ are also
compared independently in the script.

\subsection{Rebuilding}
From the package root, run
\begin{verbatim}
python -m pip install -r requirements.txt
OPENBLAS_NUM_THREADS=1 python code/verify.py --full
latexmk -pdf -interaction=nonstopmode -halt-on-error article.tex
\end{verbatim}
The environment-variable syntax is for a POSIX shell; on Windows it can
be set using the shell's equivalent syntax or omitted. A run without
\code{--full} uses smaller meshes. The numerical output records the
Python and numerical-library versions. The PDF is built from the
provided source and two generated table fragments; all are in the ZIP.

\section{Proof boundaries and a formalization route}
\label{sec:audit}

\subsection{Dependency audit}
The mathematical dependency chain is short enough to audit explicitly.
The pressure theorem requires the atomic spectral gap
(\cref{prop:gap}), the strong--weak perturbation identity
(\cref{lem:sw}), the weak pairing bound and scalar coefficient
(\cref{thm:weak}), and the pressure--moment comparison
(\cref{lem:sampling}). The atomic gap follows from an explicit weighted
Riemann sum and a finite-iterate H\"older estimate. The weak pairing
bound follows from a one-zero scaling calculation and cancellation of
its odd tail; independent zero motions are added by a partition of
unity with controlled mixed terms. No assumption of strictly positive
weights, no unproved linear-response theorem, and no
numerical spectral approximation is inserted into this chain.

The special functions only evaluate the universal integral and the
amplitude. Leaving those two quantities as their convergent integrals
would preserve the qualitative cusp and finite-size statements. The
closed forms make the laws explicit.

The article proves a local result relative to a specific family. In
particular, it does not prove a general universality theorem for all
masks with the same number of zeros. The atomic coboundary and its
spectral gap are indispensable model-specific structure. Independent
motions are allowed because they start from exactly this atomic mask.

\subsection{A staged Lean implementation}
No Lean file is included or claimed to compile. A practical future
formalization can isolate three layers. The finite layer consists of
the root-factor mask identity, the digit-product coefficient formula,
and the degree bound $b^n-1$. The real-analysis layer consists of the
one-zero integral estimates, principal-value functionals on H\"older
tests, and the weighted quadrature lemma. The operator layer consists
of rank-one exponential relaxation, resolvent projections, and
\cref{lem:sw}.

The strong--weak eigenvalue lemma is a particularly reusable target:
once a suitable spectral-projection API exists, its central equality
is just
\[
 \lambda-\beta-\nu(\Delta h)=\nu\Delta(h_\Delta-h).
\]
The essential formal issue is preserving the different norm estimates,
not replacing both by one ambient operator norm. For the singular
integrals, symmetric cutoff limits should be explicit definitions;
otherwise a principal value can be mistaken for an absolutely
convergent integral.

A useful executable certificate project would prove rigorous
finite-dimensional upper and lower spectral bounds with a separately
bounded interpolation error. The existing floating-point verifier is
not such a certificate, and a future formal checker should not treat
its output as trusted theorem data.

\section{Further research questions}
\label{sec:questions}

\begin{question}[The sharp next term]
Can the remainder in \cref{eq:main} be sharpened to an asymptotic term?
Determine whether the next correction is an ordinary quadratic term,
a fractional term, or a logarithmic or log-periodic contribution.
The present proof bounds the change of the eigenfunction only in a
strong norm and does not compute its effect at the moved zeros.
\end{question}

\begin{question}[Removing the arbitrary H\"older loss]
For fixed $0<s<1$, can one replace
$O(\delta^{1+s-\alpha})$ for every $\alpha>0$ by
$O(\delta^{1+s})$, perhaps with a logarithmic factor?
A weighted space adapted simultaneously to the atomic point and the
mask zeros could avoid choosing a smaller H\"older exponent. The
optimal total-variation law \cref{eq:TV} is a constraint on any such
argument, not by itself an answer.
\end{question}

\begin{question}[A joint limit as the moment order tends to zero]
The coefficient behaves as $J_s\sim\pi^2s^2/2$, but the proof is not
uniform as $s\downarrow0$. Determine the joint pressure asymptotic
when $s\log(1/\delta)$ remains bounded. Does the apparent quadratic
small-order response retain a universal scaling function, or do
additional orbit effects become visible?
\end{question}

\begin{question}[Independent zero motions at criticality]
At $s=1$, does the critical splitting depend only on
$\sum_j|\varepsilon_j|$, or on a different weighted combination of
zero motions? The atomic left density is no longer integrable and the
rank-one argument here fails. One must compute the perturbation of
the critical spectral block rather than substitute $s=1$ in the
subcritical coefficient.
\end{question}

\begin{question}[Uniform matching as $s\uparrow1$]
Can the compact-subcritical theorem be matched quantitatively to the
critical two-parameter law in \cite{CriticalRepo}? An overlap region
must control the blowup of $\nu_s$ and of the reduced resolvent.
The pole $J_s\sim2/(1-s)$ identifies what must be matched but does not
justify the interchange of limits.
\end{question}

\begin{question}[Two-point regularity and nonzero phases]
Is the common-phase pressure locally Lipschitz on a whole neighborhood
of the atomic phase for every fixed $0<s<1$? At nonzero phases, what
regularity is forced or obstructed by returns of singularities under
$T_b$? The estimate proved here is anchored to the atomic left
functional and does not directly give a bound between two arbitrary
nearby nonzero phases.
\end{question}

\begin{question}[Higher-multiplicity zeros]
For an atomic model with zero multiplicities $r_j>1$, what replaces
the integrability threshold and the coefficient $J_s$? A local factor
suggests replacing $s$ by $r_js$, but a pressure theorem also needs
an actual spectral gap and compatible dual density. Determine precise
hypotheses under which local multiplicities determine a tangent law.
\end{question}

\begin{question}[Response of eigenfunctions and equilibrium data]
Find a distributional or weighted-space expansion of the right and
left spectral vectors themselves. This could yield first-order laws
for observables and, in the probability-normalized subfamily,
equilibrium measures. The pressure theorem only uses a norm estimate
for the perturbed eigenfunction and does not identify its singular
profile.
\end{question}

\begin{question}[Certified computation near a cusp]
Develop interval quadrature and a proof-producing spectral enclosure
whose error remains smaller than the $O(\delta)$ pressure increment.
Uniform collocation spends many points away from the few translated
zeros. A decomposition that treats their singular parts analytically
could improve both efficiency and the trust boundary.
\end{question}

\begin{question}[An abstract moving-singularity response theorem]
Which expanding maps and finite-zero weight families admit a weak
expansion of the form ``Dirac masses plus principal values'' together
with a strong--weak eigenvalue remainder? Formulate a theorem with
explicit assumptions on the atomic spectral data, zero separation,
local exponents, and test-function regularity. The digital masks here
provide a complete worked example, not a proof for all such systems.
\end{question}

\section{Conclusion}
The lower cutoff in the previous subcritical argument was a limitation
of its perturbation estimate, not a spectral transition. Pairing the
perturbation with the atomic left eigenfunctional preserves the
cancellation around translated zeros and makes the nonlinear remainder
smaller than the cusp for every $0<s<1$.
The resulting theorem closes the repository's explicitly remaining
subcritical interval, supplies a quantitative error, and extends the
answer to independent zero motions. In this larger parameter space the
first pressure increment is $J_s\norm\eps_1$, while finite products
have the scaling amplitude $A_s$ and time scale $n\norm\eps_\infty$.
The unresolved questions now concern sharper corrections, endpoint
matching, nonzero phases, and certified implementation, rather than
whether the leading fixed-subcritical cusp exists.

\begin{thebibliography}{9}
\bibitem{GKS}
P.~Gohlke, M.~Kesseb\"ohmer, and T.~I.~Schindler,
\emph{Generalized Thue--Morse measures: spectral and fractal analysis},
arXiv:2509.22109, 2025.
\url{https://arxiv.org/abs/2509.22109}.
Theorems~2.3 and 2.7 and Section~6 describe the pressure framework and
atomic comparison relevant here.

\bibitem{GLS}
P.~Gohlke, G.~Lamprinakis, and J.~Schmeling,
\emph{On a family of singular potentials: Parameter dependence of
thermodynamic characteristics}, arXiv:2603.19001, 2026.
\url{https://arxiv.org/abs/2603.19001}.
Theorem~2.2 proves positive-order phase continuity; the discussion
following Theorem~2.3 poses the stronger-regularity question.

\bibitem{CriticalRepo}
ProveIt research corpus, research draft prepared with ChatGPT for
V.~Reshetnikov,
\emph{Critical Cusps in Digital-Product Pressure: A Jordan collision,
square-root phase scaling, and explicit subcritical response},
29 September 2026. Source in
\path{Analysis/FabiusFunction/docs/semi-formalized-research-frontiers/drafts/thue-morse/Thue_Morse_Critical_Pressure/article.tex},
repository \url{https://github.com/VladimirReshetnikov/ProveIt},
commit \nolinkurl{458ccfb80b9940220819091da107e82bb970dcca}.
The subsection explaining the $s=1/2$ cutoff and the first research
question state the gap addressed here. Unrefereed repository manuscript.

\bibitem{DLMF}
NIST Digital Library of Mathematical Functions,
Chapter~5, especially \S5.5, functional and reflection relations for
the gamma and digamma functions.
\url{https://dlmf.nist.gov/5.5}.

\bibitem{Repo}
V.~Reshetnikov, \emph{ProveIt}, repository snapshot
\nolinkurl{458ccfb80b9940220819091da107e82bb970dcca}, inspected
29 September 2026.
\url{https://github.com/VladimirReshetnikov/ProveIt}.
The Thue--Morse branch readme distinguishes the existing integer,
fractional, and critical-pressure reports and their unreviewed status.
\end{thebibliography}

\end{document}
